% Graduation Project Proposal (3--5 pages)
% !TEX TS-program = pdflatex
\documentclass[11pt,a4paper]{article}
\usepackage[margin=2.2cm]{geometry}
\usepackage[T1]{fontenc}
\usepackage{lmodern}
\usepackage{microtype}
\usepackage{amsmath}
\usepackage{booktabs}
\usepackage{enumitem}
\usepackage[hidelinks]{hyperref}
\usepackage{xurl}
\usepackage{setspace}
\usepackage{cite}
\setstretch{1.15}

% Shared numeric assumptions (must match project_inputs.yaml)
% -----------------------------------------------------------------------
% Benban site (retained as climate validation reference)
% -----------------------------------------------------------------------
\newcommand{\BenbanLat}{24.45}
\newcommand{\BenbanLon}{32.75}

% -----------------------------------------------------------------------
% PV module baseline (used in climate validation chapter)
% -----------------------------------------------------------------------
\newcommand{\Etamod}{0.215}
\newcommand{\WaterIntensity}{1.8}
\newcommand{\Nclean}{18}
\newcommand{\Fmargin}{0.25}

% -----------------------------------------------------------------------
% Economics
% -----------------------------------------------------------------------
\newcommand{\TariffUSD}{0.020}
\newcommand{\TariffEGP}{0.97}

% -----------------------------------------------------------------------
% TEC module parameters (TEC1-12706)
% -----------------------------------------------------------------------
\newcommand{\TECImax}{6.0}
\newcommand{\TECVmax}{15.4}
\newcommand{\TECQmax}{57.2}
\newcommand{\TECDTmax}{68}
\newcommand{\TECAlpha}{0.0534}
\newcommand{\TECR}{1.98}
\newcommand{\TECK}{0.524}
\newcommand{\TECCOPnom}{0.38}

% -----------------------------------------------------------------------
% Street-lamp scale (Scale A)
% -----------------------------------------------------------------------
\newcommand{\ScaleAPV}{80}
\newcommand{\ScaleATECn}{1}
\newcommand{\ScaleAWaterDay}{0.15}
\newcommand{\ScaleACleanArea}{0.35}
\newcommand{\ScaleABattWh}{150}

% -----------------------------------------------------------------------
% Residential scale (Scale R - villa rooftop, 5 kWp)
% -----------------------------------------------------------------------
\newcommand{\ScaleRPV}{5}
\newcommand{\ScaleRTECn}{6}
\newcommand{\ScaleRWaterDay}{1.5}
\newcommand{\ScaleRCleanArea}{28}
\newcommand{\ScaleRBattWh}{1200}

% -----------------------------------------------------------------------
% Commercial scale (Scale B - rooftop, 10 kWp)
% -----------------------------------------------------------------------
\newcommand{\ScaleBPV}{10}
\newcommand{\ScaleBTECn}{12}
\newcommand{\ScaleBWaterDay}{3.0}
\newcommand{\ScaleBCleanArea}{46}
\newcommand{\ScaleBBattWh}{2500}

% -----------------------------------------------------------------------
% Industrial scale (Scale C - ground-mount, 200 kWp, stretch)
% -----------------------------------------------------------------------
\newcommand{\ScaleCPV}{200}
\newcommand{\ScaleCPdcMW}{0.200}
\newcommand{\PdcMW}{200}
\newcommand{\ScaleCTECn}{80}
\newcommand{\ScaleCWaterDay}{40}
\newcommand{\ScaleCCleanArea}{930}
\newcommand{\ScaleCBattWh}{24000}

% -----------------------------------------------------------------------
% Climate / psychrometric operating point (Egyptian summer, conservative)
% -----------------------------------------------------------------------
\newcommand{\TambSummer}{38}
\newcommand{\RHmin}{25}
\newcommand{\RHdesign}{35}
\newcommand{\TdewDesign}{18.5}
\newcommand{\TcoldDesign}{14}
\newcommand{\ThotDesign}{58}
\newcommand{\DTdesign}{44}

% -----------------------------------------------------------------------
% Robot motion (commercial/industrial rail-guided)
% -----------------------------------------------------------------------
\newcommand{\RobotSpeed}{0.05}
\newcommand{\BrushWidth}{0.25}
\newcommand{\FlowNozzle}{0.008}


\title{\textbf{Project Proposal}\\[0.35em]
\large Autonomous Thermoelectric-Based Atmospheric Water Harvesting Robot\\
for Photovoltaic Panel Cleaning\\[0.25em]
\normalsize Multi-Scale Design for Egyptian New Cities}
\author{Rahma El Hamouly\\Faculty of Engineering --- Graduation Project}
\date{\today}

\begin{document}
\maketitle

\section*{Keywords}
Autonomous PV cleaning robot; thermoelectric cooler (TEC); Peltier module; atmospheric water
harvesting (AWH); photovoltaic soiling; Egypt new cities; street-lamp PV; commercial rooftop;
industrial ground-mount; prototype design.

% -----------------------------------------------------------------------
\section{Problem Statement}
% -----------------------------------------------------------------------

Photovoltaic panels in Egyptian arid environments accumulate dust at rates of 0.3--2.0\% power
loss per day, with monthly soiling losses reaching 15--40\% during the dry season without
intervention \cite{kazem2023soiling,sarver2013pvsoiling}. Wet cleaning is the most effective
remediation, restoring 90--98\% of lost transmittance \cite{mejia2021cleaning}, but it requires
a reliable supply of rinse-quality water at or near the panel array.

Egypt's new cities --- the New Administrative Capital (NAC), New Alamein, New Mansoura, and
others --- deploy PV systems at three distinct scales: (i)~\textbf{pole-mounted street-lamp
panels} (\ScaleAPV~W), installed across large urban footprints in the tens of thousands;
(ii)~\textbf{commercial rooftop arrays} (\ScaleBPV~kWp), on buildings whose sites precede
the maturation of municipal water supply; and (iii)~\textbf{industrial ground-mount plants}
(\ScaleCPV~kWp), sited at desert peripheries where piped water is unavailable
\cite{egypt2030vision,irena2023egypt}.

Current autonomous cleaning solutions depend on mains water, on-site water trucks, or a static
filled reservoir --- all of which require either infrastructure that does not yet exist or
regular human attendance. \textbf{No commercially available autonomous PV cleaner integrates
its own water source.} This project fills that gap by designing a robot that harvests cleaning
water directly from ambient air using thermoelectric cooler (TEC) modules, enabling fully
autonomous, water-independent PV panel cleaning.

% -----------------------------------------------------------------------
\section{Objectives}
% -----------------------------------------------------------------------

\begin{enumerate}[leftmargin=*]
  \item \textbf{Validate TEC-based AWH feasibility} for Egyptian climatic conditions using
  hourly psychrometric analysis of NASA POWER data at Benban Solar Park (Aswan) as the
  worst-case arid reference site.
  \item \textbf{Design TEC condensation subsystems} for three deployment scales, with
  thermodynamic sizing (COP, cold-side temperature, condensate yield, SEC) derived from first
  principles and validated against published experimental data
  \cite{eslami2021tec,shourideh2021tec,alqaed2023awg}.
  \item \textbf{Design complete autonomous robot systems} for each scale --- mechanical structure,
  motion mechanism, cleaning head, water delivery, power supply, and control architecture ---
  capable of operating without human intervention between scheduled maintenance intervals.
  \item \textbf{Develop and build a Scale~B prototype} (commercial rooftop, \ScaleBPV~kWp),
  with a complete bill of materials, fabrication protocol, and experimental validation plan.
  \item \textbf{Quantify energy and economic performance}: self-consumption fraction,
  payback period, and soiling-loss recovery value for each scale.
\end{enumerate}

% -----------------------------------------------------------------------
\section{Role of Benban in the Methodology}
% -----------------------------------------------------------------------

Benban Solar Park (\BenbanLat$^\circ$N, \BenbanLon$^\circ$E) serves as the \textbf{climate
validation reference}, not the deployment target. Its well-characterised arid
microclimate---high ambient temperature, low relative humidity, and strong diurnal
swing---represents the most thermodynamically demanding operating condition for a TEC condenser
in Egypt. Demonstrating adequate condensate yield at Benban conditions proves feasibility for
all new-city deployment sites, which generally have higher humidity (coastal or Delta locations)
and are therefore \emph{easier} cases for TEC AWH.

All prior analysis work (hourly psychrometrics, dew-point distributions, TEC COP calculations,
SEC and LCOW comparisons against VCR, desiccant, and MOF systems) is incorporated into the
project as supporting evidence and methodology. The primary deliverable is the \textbf{robot
system design and prototype}.

% -----------------------------------------------------------------------
\section{Methodology}
% -----------------------------------------------------------------------

\subsection{Climate Analysis (Benban Reference)}

Five-year hourly NASA POWER data \cite{nasa_power} are processed by
\texttt{scripts/psychrometrics\_hourly.py} to extract dew-point temperature distributions,
condensation-viable hour counts by month, and absolute humidity profiles. The conservative
summer design point is: $T_{\mathrm{amb}} = \TambSummer\,^\circ$C, RH = \RHdesign\%,
$T_{\mathrm{dp}} = \TdewDesign\,^\circ$C, harvesting window = 6~h/night.

\subsection{TEC Subsystem Sizing}

The governing equations for a single TEC1-12706 module are:
\begin{align}
  Q_c &= \alpha\,T_c\,I - \tfrac{1}{2}\,I^2 R - K\,(T_h - T_c), \label{eq:Qc} \\
  P_{\mathrm{in}} &= \alpha\,(T_h - T_c)\,I + I^2 R, \label{eq:Pin} \\
  \mathrm{COP} &= Q_c / P_{\mathrm{in}}, \label{eq:COP}
\end{align}
with parameters $\alpha = \TECAlpha$~V/K, $R = \TECR$~$\Omega$, $K = \TECK$~W/K
\cite{riffat2003peltier,eslami2021tec}. At the design operating point
($T_c = \TcoldDesign\,^\circ$C, $T_h = \ThotDesign\,^\circ$C, $I^* = 3.2$~A):
COP~$\approx$~\TECCOPnom, $P_{\mathrm{in}} \approx 22.1$~W per module.

Condensate yield per module is estimated from the humid-air mass balance across the cold plate:
\begin{equation}
  \dot{V}_w = \eta_c \,\dot{m}_{\mathrm{air}}\,(\rho_{v,\mathrm{in}} - \rho_{v,\mathrm{out}})
  \times 10^{-3} \approx 0.026\;\text{L/h}. \label{eq:Vw}
\end{equation}
Over a 6-hour summer night, one module yields $\approx$0.155~L. Module counts for each scale
follow from the daily water demand (Table~\ref{tab:scales}).

\begin{table}[htbp]
\centering
\caption{Scale summary: deployment targets and key design parameters.}
\label{tab:scales}
\begin{tabular}{llcccc}
\toprule
\textbf{Scale} & \textbf{Application} & \textbf{PV} & \textbf{Water/day} & \textbf{TEC modules} & \textbf{Battery} \\
\midrule
A & Street lamp & \ScaleAPV~W & \ScaleAWaterDay~L & \ScaleATECn & \ScaleABattWh~Wh \\
B & Commercial rooftop & \ScaleBPV~kWp & \ScaleBWaterDay~L & \ScaleBTECn & \ScaleBBattWh~Wh \\
C & Industrial & \ScaleCPV~kWp & \ScaleCWaterDay~L & \ScaleCTECn & \ScaleCBattWh~Wh \\
\bottomrule
\end{tabular}
\end{table}

\subsection{Robot System Design}

For each scale the mechanical and control design covers:
\begin{itemize}[leftmargin=*,noitemsep]
  \item \textbf{Chassis and motion:} Scale~A: single-axis rail wiper on pole frame; Scale~B:
  wheeled row-traversing platform; Scale~C: gantry rail vehicle.
  \item \textbf{Cleaning head:} microfibre brush/squeegee bar (\BrushWidth~m wide) with
  fan-spray nozzles (\FlowNozzle~L/s each).
  \item \textbf{Water treatment:} sediment filter (5~$\mu$m) $\rightarrow$ activated carbon
  $\rightarrow$ UV-C sterilisation, targeting TDS $<$ 50~mg/L \cite{beattie2012waterspots}.
  \item \textbf{Power architecture:} MPPT charge controller, LiFePO\textsubscript{4} battery,
  buck converter for TEC current control. Robot self-powered from the PV array it cleans.
  \item \textbf{Control:} ESP32-class microcontroller running a four-state machine (Harvest,
  Standby, Clean, Emergency) driven by RTC scheduling and sensor feedback.
\end{itemize}

\subsection{Prototype Build (Scale B)}

The Scale~B prototype bill of materials totals $\approx$USD~600. Build timeline: 3 weeks
(Week~1: TEC subassembly and water treatment chain; Week~2: chassis fabrication and electronics
integration; Week~3: field integration and 72-hour autonomous cycle test).

Validation metrics: condensate yield $\geq 80\%$ of model prediction; TDS $<$ 50~mg/L;
contact force $<$ 5~N; battery SoC $>$ 20\% after full cycle; transmittance restoration
$\geq 85\%$; 72-hour unattended operation.

% -----------------------------------------------------------------------
\section{Expected Deliverables}
% -----------------------------------------------------------------------

\begin{enumerate}[leftmargin=*,noitemsep]
  \item Complete technical report (\texttt{main\_report.tex} $\rightarrow$ PDF) covering all
  chapters: climate validation, TEC sizing, mechanical design (3 scales), water delivery,
  power and control, prototype plan, and results.
  \item \textbf{Scale~B working prototype} with logged sensor data from 72-hour autonomous
  test run.
  \item Engineering drawings: system architecture block diagram, robot assembly drawing (Scales
  A and B), TEC module mounting detail, water treatment schematic, control wiring diagram.
  \item Calculation notebook (\texttt{scripts/project\_calculations.py}) covering water demand,
  TEC condensate yield, power budget, and battery sizing, with auto-generated tables.
  \item Bibliography in IEEE format (\texttt{references.bib}).
\end{enumerate}

% -----------------------------------------------------------------------
\section{Risks and Mitigation}
% -----------------------------------------------------------------------

\begin{itemize}[leftmargin=*]
  \item \textbf{Low condensate yield in summer:} mitigated by running TEC during the
  longer, more humid winter nights and implementing adaptive current control based on real-time
  RH. The system can also buffer 7--14 days of production for a single cleaning event.
  \item \textbf{TEC module degradation:} Peltier modules degrade over thermal cycling. Mitigation:
  derate current to 70\% of $I_{\max}$ during operation; specify 10-year design life based on
  supplier MTBF data.
  \item \textbf{Robot derailment or brush damage:} limit normal contact force via spring-loaded
  brush mount; install reed-switch end-stops at panel frame edges.
  \item \textbf{Cost overrun:} the USD~600 BOM is dominated by the battery (30\%); LiFePO4
  prices have fallen 40\% since 2022 and continue to decline \cite{irena2023egypt}.
\end{itemize}

% -----------------------------------------------------------------------
\section{Timeline}
% -----------------------------------------------------------------------

\textbf{Week 1 (23--30 Apr 2026):} finalise and submit this proposal; complete TEC sizing and
condensate yield calculations; write Chapters 1--5 of \texttt{main\_report.tex}; begin procuring
BOM for Scale~B prototype.

\textbf{Week 2 (30 Apr -- 7 May 2026):} write Chapters 6--9 (mechanical, water, power, prototype);
fabricate TEC subassembly and water treatment chain (Scale~B); bench-test condensate yield and
water quality.

\textbf{Week 3 (7--14 May 2026):} fabricate robot chassis; integrate electronics; run 72-hour
field test; write Chapter 10 (results and discussion); full report rebuild and submission.

\bibliographystyle{ieeetr}
\bibliography{references}

\end{document}
